\subsection{Kernel-level diagnostics}\label{sec:kdiag}
The audit above compares models. To explain \emph{why} a quantum kernel does or does not help, we also
analyse the kernels themselves. Every quantum kernel studied here is noiseless, so we simulate each
feature-map state $|\phi(x)\rangle$ once and form every Gram matrix from the stored statevectors: the
fidelity kernel as $|\Psi_A\Psi_B^{\dagger}|^{2}$, and the projected kernel from each qubit's reduced
density matrix. This costs $O(N)$ circuit simulations instead of $O(N^{2})$, agrees with the circuits
used for the QSVM to within $10^{-15}$ on every entry, and reproduces the published QSVM test metrics to
four decimal places.

\paragraph{Fair surrogates.} We compare quantum and classical kernels on the identical first 2{,}000
NSL-KDD training rows, with exact Gram matrices and the same kernel SVM solver throughout. The bandwidth $\gamma$ and
$C\in\{0.1,1,10,100\}$ are selected on the training rows only, with the same grid for quantum and
classical kernels, so neither side receives more tuning than the other. We use two bandwidth grids,
because the choice matters: a narrow one of seven values from $0.003$ to $3$, and a wide one of ten
values from $0.001$ to $30$, both logarithmically spaced. Selection is either standard 5-fold
stratified cross-validation or the shift-aware variant described in Section~\ref{sec:kcv}.

\paragraph{Geometric difference and model complexity.} Following Huang et al.~\cite{huang2021power}, for
kernels normalised to $\operatorname{Tr}K=N$ we compute
\[
g(K_C\Vert K_Q)=\sqrt{\Bigl\Vert\sqrt{K_Q}\,\sqrt{K_C}\,(K_C+\lambda I)^{-2}\sqrt{K_C}\,\sqrt{K_Q}\Bigr\Vert_{\infty}},
\]
which tends to $\sqrt{\Vert\sqrt{K_Q}K_C^{-1}\sqrt{K_Q}\Vert_{\infty}}$ as $\lambda\to 0$. A small $g$
certifies that the classical kernel can learn any function the quantum kernel can learn from comparable
data, for \emph{any} labelling of the inputs, whereas a prediction advantage requires $g$ to approach
$\sqrt{N}$. We take the minimum of $g$ over each classical kernel family and bandwidth, at $N=800$ and
$\lambda=10^{-2}$. Rank-deficient kernels, such as the linear kernel on eight features, are excluded:
their regularised inverse vanishes on the null space, which makes $g$ spuriously small. We also report
the label-dependent model complexity $s_K=\Vert\sqrt{K}(K+\lambda I)^{-1}y\Vert^{2}$ with $y\in\{\pm1\}$,
which tends to $y^{\top}K^{-1}y$ as $\lambda\to 0$.

\paragraph{Concentration.} We measure the mean and variance of the off-diagonal kernel entries as the
qubit count grows from 2 to 12. Exponential concentration~\cite{thanasilp2024concentration} would drive
this variance to zero exponentially in the number of qubits.

\paragraph{Finite shots.} Each Pauli expectation $\mu$ in the projected kernel is replaced by the mean of
$M$ outcomes in $\{\pm1\}$ drawn with $P(+1)=(1+\mu)/2$, which is what hardware returns from $M$ shots in
each of the three measurement bases, and the SVM is retrained on the resulting noisy features.

\paragraph{Classical phase kernels.} Under angle encoding every qubit remains in a product state, so its
Pauli features are exactly $(\sin x_i, 0, \cos x_i)$ and the ``quantum'' projected kernel is an RBF kernel
on classical trigonometric features. We therefore also test RBF kernels on
$\varphi(x)=(\sin x_i,\cos x_i)_{i}$, and on $\varphi$ augmented with the phases that the ZZ feature
map~\cite{havlicek2019supervised} writes into the state, namely $\sin$ and $\cos$ of $2x_i$ and of
$2(\pi-x_i)(\pi-x_{i+1})$. This isolates whether an advantage requires the quantum state itself, or only
the phases it encodes.
